\chapter{Mathematical Proofs}

In this section I will provide with both definitions and examples of different types of proofs for the reader to 
gain an overview of how to use the different proving methods.

\section{Proof By Direct Argument}

Assume \(A\) to prove to prove \(B\).

\[
	A \implies B
\]

\textbf{Example I:}

For any integer \(n\), if \(n\) is even, then \( n^2 \) is even.
We will prove this theorem by direct argument.

Assume \(n\) is an even integer. Then we can write \( n = 2k \) for some integer \(k\).

Now, we compute \( n^2 \):

\[		
	n^2 = {(2k)}^2 = 4k^2 = 2(2k^2)
\]
	
This shows that \( n^2 \) is even, as it can be expressed as \( 2m \) where \( m = 2k^2 \) is an integer.
Therefore, we conclude that if \(n\) is even, then \( n^2 \) is even.

\QED

\section{Proof by Contrapositive}

\[
	A \implies B \iff \neg B \implies \neg A
\]

\textbf{Example:}

If \(n\) is an integer such that \( n^2 \) is odd, then \(n\) is odd.

We will prove this theorem by contrapositive. The contrapositive of the statement is: If \(n\) is an 
integer such that \(n\) is even, then \( n^2 \) is even.

Assume \(n\) is even. Then we can write \( n = 2k \) for some integer \(k\).

Now, we compute \( n^2 \):
	
\[
	n^2 = {(2k)}^2 = 4k^2 = 2(2k^2)
\]

This shows that \( n^2 \) is even.

Since the contrapositive statement is true, the original statement If \( n^2 \) is odd, then \(n\) is 
odd is also true.

\QED

\section{Proof By Contradiction}

Assume the negation of \(A\) to derive a false statement.

\[
	A \implies B \iff \neg A \implies F
\]

\textbf{Example I:}

If \(n\) is an integer such that \( n^2 \) is even, then \(n\) is even.

We will prove this theorem by contradiction. Assume that \(n\) is an integer such that \( n^2 \) is 
even, but \(n\) is odd. Then we can write \( n = 2k + 1 \) for some integer \(k\).

Now, we compute \( n^2 \):

\[
	n^2 = {(2k + 1)}^2 = 4k^2 + 4k + 1 = 2(2k^2 + 2k) + 1
\]
	
This shows that \( n^2 \) is odd, which contradicts our assumption that \( n^2 \) is even. Therefore, our 
assumption that \(n\) is odd must be false, and thus, \(n\) must be even.

\QED

\textbf{Example II:}

Another example is the proof that there are infinite prime numbers.

Let us define the prime numbers \(p_1, \dots, p_n\) where we assume \(p_n\) is the biggest prime number. 
If we construct a number 

\[
	q = 1 + p_1 \cdot \cdots \cdot p_n 
\]

This number is bigger than \(p_n\), but it is not divisible by any of the other prime numbers will give us always a residual of 
1. Therefore, the assumption that there is the biggest prime number is false.

\QED 

\textbf{Example II:}

Suppose, for contradiction, that \(\sqrt{2}\) is rational. Then there exist integers \(p\) and \(q\) 
with \(q\neq0\), \(\gcd(p,q) = 1\), and

\[
	\sqrt{2}=\frac{p}{q}.
\]

Squaring both sides gives

\[
	2 = \frac{p^2}{q^2} \implies p^2=2q^2.
\]

Thus, \(p^2\) is even, so \(p\) is even. Write \(p=2k\) for some integer \(k\). Then

\[
	(2k)^2=2q^2 \implies 4k^2=2q^2 \implies q^2=2k^2,
\]

so \(q^2\) is even and hence \(q\) is even. Therefore, both \(p\) and \(q\) are even, contradicting 
\(\gcd(p,q)=1\). This contradiction shows \(\sqrt{2}\) is irrational.

\QED

\section{Proof By Induction}

For all \( n \in \N \), the sum of the first \(n\) positive integers is given by:

\[
	S(n) = 1 + 2 + 3 + \cdots + n = \frac{n(n+1)}{2}
\]

We will prove this theorem by induction on \(n\).

\textbf{Base Case:} For \( n = 1 \):

\[
	S(1) = 1 = \frac{1(1+1)}{2}
\]

The base case holds.

\textbf{Inductive Step:} Assume that the statement holds for some \( n = k \), i.e., assume that:
	
\[
	S(k) = 1 + 2 + 3 + \cdots + k = \frac{k(k+1)}{2}
\]
	
We need to show that the statement holds for \( n = k + 1 \):

\[
	S(k+1) = S(k) + (k + 1)
\]
	
By the inductive hypothesis, we have:

\begin{align*}
	S(k+1) &= \frac{k(k+1)}{2} + (k + 1) \\
	&= \frac{k(k+1)}{2} + \frac{2(k + 1)}{2}\\
    &= \frac{k(k+1) + 2(k + 1)}{2} \\
	&= \frac{(k + 1)(k + 2)}{2}	
\end{align*}

Thus, the statement holds for \( n = k + 1 \).
By the principle of mathematical induction, the statement holds for all \( n \in \N \).

\QED

\section{Strong Induction}

\emph{Strong induction} is a form of mathematical induction where instead of assuming \(A(n)\) is true in our 
inductive hypothesis, we assume \(A(j)\) is true for \(1 \le j \le n\). 

The inductive step: Show \([A(1) \land A(2) \dots A(n)] \implies A(n + 1)\). 

Sometimes, we will prove specific values are true for \(1 \le j \le n\), and sometimes just assume that 
\(A(j)\) is true for \(1 \le j \le n\). 

\textbf{Example:}

Prove that any integer \(\ge 1\) can be written as the product primes.

Let \(A(n)\) be that if \(n\) is an integer greater than 1, then \(n\) can be written as the product of 
primes. 

\textbf{Basis: } \(A(2)\) is true because is prime 

\textbf{Inductive hypothesis:} \(A(j)\) for \(2 \le j \le n\) is true. \textbf{Important!}

Consider the cases \(k + 1\) is prime or composite. 

The first case is true.

For the case \(n + 1\) is not prime then \(n + 1 = ab\) with \(2 \le a,b \le n < n + 1\) and both 
for \(a\) and \(b\) we know that they can be written as the product of primes thus 
\(n + 1\) can be written as the product of primes. 

\QED

\section{Proof By Exhaustion}

In \emph{exhaustion} we test all possible cases and their respective branches.

\textbf{Example:}

The only integer solutions to the equation \( x^2 + y^2 = 1 \) are \( (0, 1), (1, 0), (0, -1), (-1, 0) \).

We will prove this theorem by exhaustion. We will check all possible integer values of \(x\) and 
\(y\) such that \( x^2 + y^2 = 1 \).

The possible integer values for \(x\) and \(y\) are \( -1, 0, 1 \). We will check each case:

If \( x = 0 \):

Then \( y^2 = 1 \) gives \( y = 1 \) or \( y = -1 \).

Solutions: \( (0, 1), (0, -1) \).

If \( x = 1 \):

Then \( y^2 = 0 \) gives \( y = 0 \).

Solution: \( (1, 0) \).

If \( x = -1 \):

Then \( y^2 = 0 \) gives \( y = 0 \).

Solution: \( (-1, 0) \).

Thus, the only integer solutions to the equation are:
	
\[
	(0, 1), (1, 0), (0, -1), (-1, 0)
\]

\QED

\section{Proof By Counter-example}

In a \emph{counter-example} is it enough to find one case for a statement to not be true, to prove an 
statement false, and in some cases that the opposite is true.

\textbf{Example:}

The statement all prime numbers are odd is false.

To prove this theorem, we will provide a counterexample.

The number 2 is a prime number, as its only divisors are 1 and 2. However, 2 is 
even, which contradicts the statement that all prime numbers are odd.

Therefore, the statement All prime numbers are odd is false.

\QED